\documentclass[11pt]{beamer}

\usepackage[utf8]{inputenc}
\usepackage[T1]{fontenc}
\usepackage[default]{opensans}
\usepackage[spanish]{babel}
\usepackage{hyperref}
\usepackage{tikz}
\usepackage[dvipsnames]{xcolor}
\usepackage{graphicx}
\usetikzlibrary{calc,patterns}
\usepackage{bm}

\hypersetup{
  colorlinks=true,
  urlcolor=blue,
  linkcolor=blue,
  citecolor=blue
}

\usetheme{default}
\usecolortheme{dolphin}
\useinnertheme{circles}

\title{Navegación de drones consciente de la creencia bajo observabilidad parcial}
\subtitle{Propuesta de proyecto}
\author[M. Osipova]{Mariia Osipova\\[3pt]\textit{mosipova@udesa.edu.ar}}
\institute{\\[12pt]Bajo la supervisión de Juan Felipe Perdomo}
\date[Abr 2026]{Abril 2026}

%----------------------------------------------------------------------------------------
% fuentes
%----------------------------------------------------------------------------------------
\newcommand{\srcSuttonBarto}{\href{http://incompleteideas.net/book/the-book.html}{Sutton \& Barto (2018)}}
\newcommand{\srcKochenderferADM}{\href{https://algorithmsbook.com}{Kochenderfer, Wheeler \& Wray (2022)}}
\newcommand{\srcKaelbling}{Kaelbling, Littman \& Cassandra (1998)}
\newcommand{\srcSchulman}{\href{https://arxiv.org/abs/1707.06347}{Schulman et al. (2017)}}
\newcommand{\srcTobin}{Tobin et al. (2017)}
\newcommand{\srcHwangbo}{Hwangbo et al. (2017)}
\newcommand{\SourceLine}[1]{\vfill{\tiny\textcolor{gray}{Fuentes: #1}}}

%----------------------------------------------------------------------------------------

\begin{document}

%----------------------------------------------------------------------------------------
% PORTADA
%----------------------------------------------------------------------------------------

\begin{frame}
    \titlepage
\end{frame}

%----------------------------------------------------------------------------------------
% ÍNDICE
%----------------------------------------------------------------------------------------

\begin{frame}{Resumen de la presentación}
  \tableofcontents
\end{frame}

%------------------------------------------------

\section{Motivación}

%------------------------------------------------

\begin{frame}{¿De qué se trata el proyecto?}

En términos simples: comparar tres maneras de controlar un dron simulado cuando los sensores no son perfectos.

\medskip
\textbf{Tarea base:} volar desde el origen hasta un destino fijo.

\medskip
Los tres enfoques a comparar:
\begin{itemize}
  \item \textbf{PID} --- un controlador clásico escrito a mano,
  \item \textbf{RL sin información extra} --- un agente que aprende solo con lecturas ruidosas,
  \item \textbf{RL con ``creencia''} --- el mismo agente, pero que además sabe cuánto confía en sus sensores.
\end{itemize}

\medskip
\begin{block}{Pregunta central}
¿Ayuda que el agente sepa cuánto no sabe?
\end{block}

\smallskip
{\small Todo en simulación --- sin dron físico.}
\end{frame}

%------------------------------------------------

\begin{frame}{¿Por qué este problema?}

La pregunta detrás de este proyecto es simple:

\medskip
\begin{quote}
¿Cómo debería actuar un dron cuando no sabe exactamente dónde está?
\end{quote}

\medskip
Lo que volvió atractiva esta pregunta para mí fue la combinación de:
\begin{itemize}
  \item aprendizaje por refuerzo como forma de aprender comportamiento por interacción,
  \item robótica como banco de pruebas concreto para la toma de decisiones,
  \item incertidumbre como la condición que vuelve al problema genuinamente difícil.
\end{itemize}

El contexto es la navegación autónoma de un cuadricóptero con sensado ruidoso e incompleto.

\SourceLine{\srcSuttonBarto; \srcKochenderferADM}
\end{frame}

%------------------------------------------------

\begin{frame}{Del control a la observabilidad parcial}

Mi camino hacia este tema fue gradual:
\begin{itemize}
  \item primero, control clásico: por ejemplo, un controlador PID para estabilización de altitud;
  \item luego, aprendizaje por refuerzo: dejar que el agente aprenda en lugar de codificar el comportamiento a mano;
  \item finalmente, observabilidad parcial: un vuelo real implica datos de sensores ruidosos, retrasados e incompletos.
\end{itemize}

\medskip
Eso cambia la naturaleza de la tarea:
\begin{block}{Cambio clave}
El dron no debería actuar solo sobre observaciones crudas. Debería actuar sobre lo que cree acerca de su estado, junto con cuán seguro está de esa creencia.
\end{block}

\SourceLine{\srcSuttonBarto; \srcKaelbling}
\end{frame}

%------------------------------------------------

\section{Formulación del problema}

%------------------------------------------------

\begin{frame}{MDP vs. POMDP}

En una formulación estándar de RL, el agente observa el estado verdadero y aprende una política
\[
\pi(a \mid s).
\]

Para un dron con sensores imperfectos, ese supuesto es falso.

\medskip
Un modelo más adecuado es un POMDP:
\begin{itemize}
  \item el estado verdadero está oculto,
  \item el agente recibe observaciones ruidosas $o_t$,
  \item el objeto relevante es una creencia sobre los posibles estados.
\end{itemize}

\medskip
Entonces el problema central no es solo \emph{qué acción tomar}, sino también \emph{qué creer sobre el estado actual}.

\SourceLine{\srcKochenderferADM; \srcKaelbling}
\end{frame}

%------------------------------------------------

\begin{frame}{Pregunta de investigación}

\begin{block}{Pregunta central}
¿Puede un agente de aprendizaje por refuerzo navegar un cuadricóptero de manera más efectiva bajo datos de sensores ruidosos e incompletos al mantener explícitamente una creencia sobre su propio estado, en comparación con un agente que trata las observaciones crudas como verdad de terreno?
\end{block}

\medskip
Más específicamente, quiero evaluar si la representación explícita de la incertidumbre mejora:
\begin{itemize}
  \item la robustez frente al ruido,
  \item la estabilidad durante el entrenamiento,
  \item la eficiencia en muestras,
  \item la generalización a condiciones no vistas.
\end{itemize}
\end{frame}

%------------------------------------------------

\section{Método}

%------------------------------------------------

\begin{frame}{Configuración experimental}

El proyecto usará una tarea de navegación de cuadricóptero simulado construida con:
\begin{itemize}
  \item \texttt{gym-pybullet-drones} para el entorno,
  \item Stable-Baselines3 para PPO,
  \item NumPy / SciPy para estadísticas de creencia,
  \item Matplotlib para trayectorias y curvas de entrenamiento.
\end{itemize}

\medskip
Tarea base:
\begin{itemize}
  \item volar desde el origen hasta un waypoint fijo.
\end{itemize}

Controladores a comparar:
\begin{itemize}
  \item PID,
  \item PPO con observaciones crudas,
  \item PPO con observaciones aumentadas con creencia.
\end{itemize}

\SourceLine{\srcSchulman; \srcHwangbo}
\end{frame}

%------------------------------------------------

\begin{frame}{Fase 1: línea de base bajo condiciones ideales}

Objetivo: establecer si la canalización de RL funciona antes de introducir incertidumbre.

\medskip
Voy a comparar:
\begin{itemize}
  \item un controlador PID clásico con acceso directo al estado,
  \item un agente PPO entrenado desde cero en la misma tarea de waypoint.
\end{itemize}

Resultados principales de esta fase:
\begin{itemize}
  \item curvas de recompensa,
  \item trayectorias exitosas,
  \item una línea de base limpia para estudios posteriores de degradación.
\end{itemize}

\SourceLine{\srcSchulman}
\end{frame}

%------------------------------------------------

\begin{frame}{Fase 2: introducir incertidumbre}

Voy a degradar la canalización de observación usando tres mecanismos de ruido:
\begin{itemize}
  \item ruido gaussiano sobre la posición observada,
  \item pérdida de observaciones o repetición de mediciones viejas,
  \item perturbaciones externas de viento.
\end{itemize}

\medskip
Esta fase está diseñada para producir curvas de degradación y responder:
\begin{quote}
¿Qué tan rápido colapsa el desempeño cuando el sensado se vuelve menos confiable?
\end{quote}

\medskip
También prueba si la política RL simple aprendió robustez genuina o si simplemente se sobreajustó a observaciones limpias.

\SourceLine{\srcKochenderferADM; \srcSchulman}
\end{frame}

%------------------------------------------------

\begin{frame}{Fase 3: observaciones aumentadas con creencia}

Esta es la contribución central.

\medskip
A partir de una ventana deslizante de observaciones recientes, calcular:
\[
\hat{\mu}_t \qquad \text{y} \qquad \hat{\sigma}_t,
\]

donde:
\begin{itemize}
  \item $\hat{\mu}_t$ es una estimación filtrada del estado,
  \item $\hat{\sigma}_t$ es una medida de confianza.
\end{itemize}

La política recibe entonces una entrada aumentada:
\[
[o_t, \hat{\mu}_t, \hat{\sigma}_t].
\]

Hipótesis: el agente aprenderá a comportarse con más cautela cuando la incertidumbre sea alta y con mayor decisión cuando sea baja.

\SourceLine{\srcKaelbling; \srcSchulman}
\end{frame}

%------------------------------------------------

\begin{frame}{Fase 4: generalización y transferencia}

Voy a comparar tres regímenes de entrenamiento:
\begin{itemize}
  \item un especialista entrenado solo con ruido bajo,
  \item un especialista entrenado solo con ruido alto,
  \item un generalista entrenado con aleatorización del dominio a través de niveles de ruido.
\end{itemize}

\medskip
La evaluación se hará sobre condiciones de ruido no vistas.

\medskip
Pregunta clave:
\begin{block}{Pregunta de transferencia}
¿El generalista aumentado con creencia transfiere mejor a condiciones no vistas que las políticas especializadas en un solo régimen?
\end{block}

\SourceLine{\srcTobin}
\end{frame}

%------------------------------------------------

\section{Resultados esperados}

%------------------------------------------------

\begin{frame}{Resultados esperados}

\begin{itemize}
  \item Bajo condiciones ideales, PPO debería igualar aproximadamente a PID después de suficiente entrenamiento.
  \item PID debería ser más eficiente en muestras porque requiere ajuste, no aprendizaje.
  \item Bajo ruido moderado y alto, PPO simple podría degradarse bruscamente si se sobreajusta a observaciones limpias.
  \item PPO aumentado con creencia debería superar a PPO simple cuando el sensado se vuelve poco confiable.
  \item El entrenamiento con aleatorización del dominio debería mejorar la transferencia a regímenes de ruido no vistos.
\end{itemize}

\medskip
El resultado principal esperado es que incluso una señal de creencia liviana puede mejorar el comportamiento de RL bajo observabilidad parcial.
\end{frame}

%------------------------------------------------

\begin{frame}{Productos y evaluación}

Productos concretos:
\begin{itemize}
  \item curvas de entrenamiento para todas las variantes de agentes,
  \item gráficos de trayectorias para PID, PPO simple y PPO consciente de la creencia,
  \item curvas de degradación RMSE bajo distintos niveles de ruido,
  \item una matriz de generalización train/test a través de regímenes de ruido,
  \item un encuadre escrito de la tarea como POMDP.
\end{itemize}

\medskip
Criterios de evaluación:
\begin{itemize}
  \item precisión de navegación,
  \item tasa de éxito,
  \item robustez frente a la incertidumbre,
  \item generalización fuera del régimen de entrenamiento.
\end{itemize}
\end{frame}

%------------------------------------------------

\section{Discusión}

%------------------------------------------------

\begin{frame}{Por qué importa}

Este proyecto es una prueba pequeña pero concreta de una pregunta más amplia en aprendizaje por refuerzo:

\medskip
\begin{quote}
¿Qué representaciones deben hacerse explícitas para que un agente actúe bien bajo incertidumbre?
\end{quote}

\medskip
Si las entradas conscientes de la creencia ayudan, eso respaldaría dos afirmaciones:
\begin{itemize}
  \item la observabilidad parcial es un cuello de botella práctico, no solo un detalle teórico,
  \item la información explícita sobre incertidumbre puede cambiar el comportamiento que aprende un agente de RL.
\end{itemize}

También conecta control, robótica, RL y razonamiento basado en POMDP dentro de un único entorno experimental concreto.

\SourceLine{\srcSuttonBarto; \srcKochenderferADM; \srcKaelbling}
\end{frame}

%------------------------------------------------

\begin{frame}{Limitaciones y extensiones}

Limitaciones actuales:
\begin{itemize}
  \item la representación de creencia es solo una aproximación liviana,
  \item PPO es solo un algoritmo posible de RL,
  \item el proyecto sigue estando en simulación.
\end{itemize}

\medskip
Posibles extensiones:
\begin{itemize}
  \item comparar con estimación del estado basada en filtro de Kalman,
  \item probar políticas recurrentes como LSTMs,
  \item explorar RL basado en modelo o solucionadores explícitos de POMDP,
  \item estudiar más adelante la transferencia sim-to-real.
\end{itemize}
\end{frame}

%------------------------------------------------

\section{Referencias}

%------------------------------------------------

\begin{frame}{Referencias seleccionadas}
\footnotesize
\begin{itemize}
  \item Sutton, R. S., \& Barto, A. G. (2018). \textit{Reinforcement Learning: An Introduction}.
  \item Kochenderfer, M. J., Wheeler, T. A., \& Wray, K. H. (2022). \textit{Algorithms for Decision Making}.
  \item Kaelbling, L. P., Littman, M. L., \& Cassandra, A. R. (1998). Planning and acting in partially observable stochastic domains.
  \item Schulman, J. et al. (2017). Proximal Policy Optimization Algorithms.
  \item Hwangbo, J. et al. (2017). Control of a quadrotor with reinforcement learning.
  \item Tobin, J. et al. (2017). Domain randomization for transferring deep neural networks from simulation to real world.
\end{itemize}
\end{frame}

%------------------------------------------------

\begin{frame}
  \centering
  \Huge Gracias
\end{frame}

\end{document}
